\hspace{3mm}

\centerline{\bf \huge 10 класс}
\centerline{\normalsize{\it Работа рассчитана на \bf{210} минут}}

\hspace{3mm}

\problem{10.1} Найдите значение выражения: \\

$\Biggl(1 + 2\sqrt{506}\Biggl) \Biggl( \dfrac{1}{\sqrt{1} + \sqrt{2}}+\dfrac{1}{\sqrt{2} + \sqrt{3}}+\dotsc+\dfrac{1}{\sqrt{2023} + \sqrt{2024}}\Biggr)$

\hspace{3mm}

Ответ: 2023.

Решение:\\

$\Biggl(1 + 2\sqrt{506}\Biggl) \
\Biggl( \dfrac{1}{\sqrt{1} + \sqrt{2}}\cdot\dfrac{\sqrt{2}-\sqrt{1}}{\sqrt{2} - \sqrt{1}}+
\dfrac{1}{\sqrt{2} + \sqrt{3}}\cdot\dfrac{\sqrt{3}-\sqrt{2}}{\sqrt{3} - \sqrt{2}} + \dotsc +\\
\dfrac{1}{\sqrt{2023} + \sqrt{2024}}\cdot\dfrac{\sqrt{2024}-\sqrt{2023}}{\sqrt{2024} - \sqrt{2023}}\Biggr) = 
\Biggl(1 + 2\sqrt{506}\Biggl)\Biggl( \dfrac{\sqrt{2} - \sqrt{1}}{2 - 1}+\dfrac{\sqrt{3} - \sqrt{2}}{3-2}+\dotsc+\dfrac{\sqrt{2024} - \sqrt{2023}}{2024 - 2023}\Biggr) = (1 + 2\sqrt{506})(\sqrt{2} - \sqrt{1} + \sqrt{3} - \sqrt{2} + \dotsc + \sqrt{2023} - \sqrt{2022} + \sqrt{2024} - \sqrt{2023}) = (1 + 2\sqrt{506})(-1 + \sqrt{2024}) = (\sqrt{4\cdot 506}+1)(\sqrt{2024} - 1) = (\sqrt{2024} +1 )(\sqrt{2024} - 1) = 2024 - 1 = 2023$

\hspace{3mm}

Критерии:\\
Приведено полное обоснованное решение: \textbf{7 баллов.}\\
Идеи, продвигающие ход решения, но недостаточные для получения верного ответа: \textbf{2 балла.}\\
Неверное решение или идеи, не продвигающие ход решения: \textbf{0 баллов.}\\
Верный или неверный ответ без обоснования: \textbf{0 баллов.}

\hspace{3mm}

\problem{10.2}В 9:00 19 января $2023$ года на доске были записаны два числа: $2023$ слева и $2024$ справа. Каждую минуту, начиная с 9:00, Арслан будет стирать числа с доски и вместо них записывать слева число $\dfrac{2023ab}{a+b}$, а справа число $\dfrac{1}{a} + \dfrac{1}{b}$ (где a - левое, а b - правое число на доске до стирания). Чему будет равно произведение чисел на доске в $23:59 \\ 19$ января $2023$ года?

\hspace{3mm}

Ответ: $2023$.

\newpage

\hspace{3mm}

Решение:\\
Заметим, что произведение чисел на доске, начиная с 9:01, не изменяется и всегда равно $2023$. 
Пусть в какой-то момент времени мы стерли два числа a и b и записали вместо них: $\dfrac{2023ab}{a+b} \ \text{и} \ \dfrac{1}{a} + \dfrac{1}{b},$ произведение которых равно $\dfrac{2023ab}{a+b} \cdot \Biggl( \dfrac{1}{a} + \dfrac{1}{b} \Biggr) = \dfrac{2023ab}{a+b} \cdot \dfrac{a+b}{ab} = 2023.$
Значит, и в $23:59$ 19 января $2023$ года произведение будет равно $2023$.

\hspace{3mm}

Критерии:\\
Приведено полное обоснованное решение: \textbf{7 баллов.}\\
Идеи, продвигающие ход решения, но недостаточные для получения верного ответа: \textbf{2 балла.}\\
Неверное решение или идеи, не продвигающие ход решения: \textbf{0 баллов.}\\
Верный или неверный ответ без обоснования: \textbf{0 баллов.}


\hspace{3mm}

\problem{10.3}Джангар каждый будний день недели побеждает некоторое количество противников, а в выходные дни он отдыхает. Известно, что в каждый последующий день недели он побеждал не больше противников, чем в предыдущий. И за всю третью неделю декабря, начиная с понедельника, он победил 6000 врагов. Какое наименьшее суммарное количество врагов он мог победить за первый, третий и пятый дни в третью неделю декабря?

\hspace{3mm}


Ответ: 3000 врагов.

Решение:\\
\textit{Оценка:} Докажем, что в первый, третий и пятый день меньше 3000 побеждённых врагов быть не могло. 
Действительно,  пусть в эти дни побеждено меньше 3000 врагов. Так как во второй и в четвёртый день богатырь победил не больше, чем в первый и в третий, то во второй и в четвёртый также побеждено меньше 3000 врагов. Тем самым в сумме за пять дней побеждено меньше 6000 врагов, что противоречит условию. 

\textit{Пример:} Если в каждый из первых четырёх дней богатырь побеждал по 1500 врагов, а в пятый - никого, то условия задачи выполнены. И за первый, третий и пятый дни он победил ровно 3000 врагов.

\newpage

\hspace{3mm}

Критерии:\\
Приведена и доказана верная оценка на 3000 врагов и приведён любой правильный пример на 3000 врагов: \textbf{7 баллов.}\\
Приведена и доказана верная оценка на 3000 врагов, но пример не приведён или приведён неправильный: \textbf{5 баллов.}\\
Приведён любой правильный пример - 3000 врагов, но не доказано, почему не может быть меньше:\textbf{2 балла.}\\
Приведена неверная оценка или приведён неверный пример: \textbf{0 баллов.}
Верный или неверный ответ без обоснования: \textbf{0 баллов.} 

\hspace{3mm}

\problem{10.4}Известно, что для любого $x$ выполнено неравенство:\\ $ax^{20} + bx^{12} + 202c \geq bx^{20} + cx^{12} + 2023a \geq cx^{20} + ax^{12} + 2023b$.\\ Докажите, что $a = b = c.$

\hspace{3mm}

Решение:\\
Из неравенства имеем: \\
$(ax^{20} + bx^{12} + 2023c) - (bx^{20} + cx^{12} + 2023a) \geq 0$, следовательно, $(a-b)x^{20} + (b-c)x^{12} + 2023(c-a) \geq 0.$ Так как трёхчлен не отрицателен при любом значении x, то старший коэффициент $(a-b) \geq 0$\\
$(bx^{20} + cx^{12} + 2023a) - (cx^{20} + ax^{12} + 2023b) \geq 0$, следовательно, $(b-c)x^{20} + (c-a)x^{12} + 2023(a-b) \geq 0.$ Так как трёхчлен не отрицателен при любом значении $x$, то старший коэффициент $(b-c) \geq 0$\\
А это значит, что $a \geq b$ и $b \geq c$ перепишем в виде $a \geq b \geq c$

Так как исходные неравенства выполнены при любых x, то подставим $x = 0$
Имеем:
$2023c \geq 2023a \geq 2023b$ , а значит $c \geq a \geq b$

Получается, $a \geq b \geq c$ и $c \geq a \geq b$, следовательно $c \geq a \geq b \geq c$, а данное неравенство имеет смысл только при $a = b = c.$

\hspace{3mm}

Критерии:\\
Приведено полное обоснованное решение: \textbf{7 баллов.}\\
Доказано только, что $a \geq b \geq с$: \textbf{4 балла.}\\
Доказано только, что $c \geq a \geq b$: \textbf{3 балла.}\\

\newpage

\hspace{3mm}

\noindentИдеи, продвигающие ход решения, но недостаточные для получения верного ответа: \textbf{1 балл.}\\
Неверное решение или идеи, не продвигающие ход решения: \textbf{0 баллов.}

\hspace{3mm}

\problem{10.5} В треугольнике $ABC$ точки $M$ и $N$ — середины сторон $AC$ и $BC$ соответственно. Известно,
что точка пересечения медиан треугольника $AMN$ является точкой пересечения высот треугольника
$ABC$. Найдите угол $ABC$.

\hspace{3mm}

Ответ: $45^{\circ}$

Решение:\\
Пусть $H$ – ортоцентр (точка пересечения высот) треугольника $ABC$. Тогда высота $AT$ треугольника $ABC$ содержит медиану треугольника $AMN$, то есть пересекает отрезок $MN$ в его середине — точке $E$ (см. рис.).

\begin{figure}[h] 
    \includegraphics[scale=0.7]{Городская Олимпиада Элиста 2023/Решения/picture_10.5_sol.jpg}
\end{figure}

Так как $MN \parallel AB$, то треугольники $ETN$ и $ATB$ подобны, следовательно $\dfrac{TN}{TB} = \dfrac{TE}{TA} = \dfrac{EN}{AB} = \dfrac{1}{4}.$ Пусть $TN = x, TE = y.$ Тогда $BN = 3x, AE = 3y.$ Следовательно, $CT = CN - TN = 2x,$ а $EH = \dfrac{1}{3} AE = y$ (по свойству точки пересечения медиан треугольника). 

Заметим, что в прямоугольных треугольниках $CTH$ и $BTA \ \text{имеем} \ \dfrac{CT}{BT} = \dfrac{2x}{4x} = \dfrac{1}{2}$ и $\dfrac{HT}{AT} = \dfrac{2y}{4y} = \dfrac{1}{2}.$ Значит треугольники подобны. Следовательно $\angle TCH = \angle TBA.$ Но $CH - $ часть высоты $CQ$ треугольника $ABC$, поэтому эти равные углы являются острыми углами прямоугольного треугольника $CQB$, то есть каждый из них равен $45^{\circ}.$

\newpage

\hspace{3mm}

Критерии:\\
Приведено полное обоснованное решение: \textbf{7 баллов.}\\
Приведено верное в целом рассуждение, в котором допущены несущественные пробелы или неточности: \textbf{5 баллов.}

\newpage

\hspace{3mm}

\noindentИдеи, продвигающие ход решения, но недостаточные для получения верного ответа: \textbf{2 балла.}\\
Неверное решение или идеи, не продвигающие ход решения: \textbf{0 баллов.}\\
Верный или неверный ответ без обоснования: \textbf{0 баллов.}

